\newcommand{\figImpulseFuncGraph}{
    \begin{figure}[htbp]
    \centering
    \begin{tikzpicture}[scale=1.5]
    % 坐标轴
    \draw[->] (-3,0) -- (3,0) node[right] {$t$};
    \draw[->] (0,-0.5) -- (0,2) node[above] {$\delta(t)$};
    % 冲激箭头
    \draw[->, thick, red] (0,0) -- (0,1.5) node[above right] {$\infty$};
    % 原点标记
    \fill (0,0) circle (2pt);
    \node[below left] at (0,0) {$0$};
    % 积分面积示意
    \draw[blue, dashed] (-0.3,0) -- (-0.3,0.5) -- (0.3,0.5) -- (0.3,0);
    \node[blue] at (0,0.7) {$1$};
    \end{tikzpicture}
    \caption{冲激信号函数图像}
    \label{fig:ImpulseFuncGraph}
    \end{figure}
}


\newcommand{\figResistanceUIAnalysis}{
    \begin{figure}[htbp]
    \centering
    \begin{circuitikz}

    \draw (0,0) to [short, o-] (2,0)
                to[R, l = $R$] (4,0)
                to [short, -o] (6,0);
    \node[below] at (0,0) {$+$};
    \node[below] at (6,0) {$-$};
    \node[below] at (3,-1) {$U$};
    
    \draw[->] (2, 1) -- (4, 1);
    \node[above] at (3,1) {$I$};

    \end{circuitikz}
    \caption{电阻电路}
    \label{fig:ResistanceUIAnalysis}
    \end{figure}
}


\newcommand{\figResistanceUIAnalysisFancImage}{
    \begin{figure}[htbp]
    \centering
    \begin{tikzpicture}[scale=0.8]
        % 坐标轴
        \draw[->] (-0.5,0) -- (5,0) node[right] {$u$};
        \draw[->] (0,-0.5) -- (0,4) node[above] {$i$};
        % 绘制电阻UI曲线
        \draw[domain=0:3, smooth, thick] plot (\x,{0.8*\x});
        % 原点标记
        \fill (0,0) circle (2pt);
        \node[below left] at (0,0) {$O$};
    \end{tikzpicture}
    \caption{电阻UI曲线}
    \label{fig:ResistanceUIAnalysisFancImage}
    \end{figure}
}
